%/*% pdflatex -file-line-error -interaction nonstopmode %

\documentclass{article}
\usepackage{amsmath}
\usepackage{amsfonts}
\usepackage{hyperref}
\usepackage{exercises}
\usepackage{amsthm}
\usepackage{yhmath}
\newcommand*{\oldneg}{\mathord{\sim}}

\theoremstyle{definition}
\newtheorem{definition}{Definition}

\setlength{\parindent}{0pt}
\setlength{\parskip}{10pt}

\title{Chapter 12 Answers}
\begin{document}
\maketitle

\section*{Exercise}
$Note$: Smullyan omits saying what alphabet we should show $Ax$ represents the set of all alternating strings, so I will assume he meant the alphabet 
$\{a,b\}$ (as this is the one he defines (E) with).\newline
We are to show that $Ax$ really represents the set of all alternating strings over the alphabet $\{a,b\}$.
\begin{proof} On the one hand, assume that $Ax$ is provable. In the base case where $x$ $a,b,ab$ or $ba$, these are all alternating by definition. 
So let us make an induction assumption that for all $Ax$ provable with $x$ of length $n$, $x$ is alternating. Well, the only formulas of length $n+1$ 
for which $Ax$ could be provable would be of the form $Axab$ where $Axa$ is provable, or of the form $Axba$ where $Axb$ is provable. In both cases, we 
have $Axa$, $Axb$ having $xa$ and $xb$ both being alternating by the assumption, thus meaning it is obvious that $Axab, Axba$ would both be alternating, 
with respect to the provable sentences $xa, xb$.\newline \newline
Conversley, we wish to show that all alternating strings $x$ in the alphabet $\{a,b\}$ are such that $Ax$ is provable. The base cases of lengths $1$ and $2$ 
are given by the first four facts of the system. Now, make an induction assumption that all alternating strings $s$ from the alphabet ${a,b}$ of degree $n$ 
are such that $Ax$ is provable. Thus $s$ must either be of the form $xa$ or $xb$ for some $x$. Hence, because $s$ is either $xa$ or $xb$, 
either $Axab$ or $Axba$ 
are provable in $(E)$, respectively. Also, the only ay to make an alternating string starting from some string of the from $xa$ or $xb$ is from the string 
$xab$ or $xba$. Hence, all of the alternating strings of degree $n+1$ are such that $Ax$ is provable, hence proving that all alternating strings 
from the alphabet $\{a,b\}$ are such that $Ax$ is provable, thus completing the proof.
\end{proof}
\section*{Problem 1}
We are considering two $EFS$'s $E_1, E_2$ over the alphabet $K$ in which the sets $W_1, W_2$ are represented, respectively. Well, we consider all of 
those predicates common to both systems. We then replace these common predicates in $E_2$ by other symbols of the predicates, and keep all of the predicates 
in $E_1$. Then, we consider the $EFS$ $E$ which is the combination of both of the axiom schemes of $E_1$ and the altered system of $E_2$. Because all 
the axiom schemes are the same, the only thing that could have made $W_1, W_2$ no longer representable is if 'extra' axiom schemes using the same 
predicates as the ones that represent them were added. However, this was guaranteed not to happen as all of the possible 'added' predicates 
would be due to predicates in common to $E_1, E_2$, which we changed when we redefined $E_2$. Thus, $W_1$ and $W_2$ are still both representable.
\section*{Problem 2}
We suppose $W_1, W_2$ are representable over $K$. Hence, by Problem 1, there is some $(E)$ over $K$ such that both $W_1, W_2$ are representable. So, 
we simply add to $(E)$ the predicate $P$ with the following axioms:\newline $P_1x\rightarrow Px$\newline $P_2x\rightarrow Px$ for $P_1, P_2$ the predicates 
that represent $W_1, W_2$, respectively. Clearly, $P$ represents $W_1\bigcup W_2$. On the other hand, we add the axiom:\newline
$P_1x\rightarrow P_2x\rightarrow Px$ meaning $Px$ represents $W_1\bigcap W_2$.
\section*{Problem 3}
We assume $W_1, W_2$ are both solvable over $K$. Hence, all of $W_1, \widetilde W_1, W_2, \widetilde W_2$ are representable over $K$. Now, by Problem 2, 
$W_1\bigcup W_2$ and $W_1\bigcap W_2$ are representable over $K$. Also, $\widetilde W_1\bigcap \widetilde W_2$ and $\widetilde W_1\bigcup \widetilde W_2$ 
are both representable over $K$, that is, both $\widetilde {W_1\bigcup W_2}$ and $\widetilde {W_1\bigcap W_2}$ are representable over $K$, thus proving
 that both $W_1\bigcup W_2$ and $W_1\bigcap W_2$ are solvable over $K$.
\section*{Problem 4}
(a):\newline This problem really asks us to find what it means for $x$ to be the successor of $y$ in dyadic notation. Well, the base cases are simple: 
$S 1,2; S 2,11$. Then, it is obvious that the successor of some dyadic representation of the form $x1$ is $x2$. Lastly, we prove that the successor of 
some $x2$ is $y1$ for $y$ the successor of $x$. Well, if $x2$ is some dyadic representation, and $y$ is the successor of $x$, then $y1$ is just the 
concatenation of $y$ and $1$, that is, $x+1$ and $1$, which is equal to $(x+1)\times 2^1+1=2x+3$, while $x2$ is $2x+2$, thus proving $S x2, y1$ for 
$S x,y$. Hence, the $EFS$ over $D$ with the axioms schemes:\newline
$S1,2$\newline $S2,11$\newline $Sx1,x2$\newline $Sx,y\rightarrow Sx2,y1$\newline is such that $Sx,y$ dyadically enumerates the relation $x$ is the successor 
of $y$
(as can easily be shown by induction).
\newline (b):\newline Consider the following $EFS$ over $D$ (and $Lx,y$ is to be read as $x$ is less than $y$):\newline 
$Sx,y\rightarrow Lx,y$\newline $Lx,y\rightarrow Ly,z\rightarrow Lx,z$\newline
Well, it is obvious that any $Lx,y$ provable in the system is such that $x$ is really less than $y$. Conversley, assume that $x<y$. Thus $x<z$ for $Sx,z$.
Then, we can use the same rule to know that $Lz,v$ for $Sv,z$, and thus we can conclude that $Lx,v$. And we can then repeat this process on $v$, and so on
(that is, we show $Sv,l\rightarrow Lx,l$ and, in general, if we have some $Lx,n$, we use the second axiom scheme to conclude that $Lx,k$ for $Sn,k$),
until we finally reach $y$ (and we will at some point), which lets us prove $Lx,y$, thus completing the proof.\newline
(c):\newline Consider the following axiom scheme:\newline
$Sx,y\rightarrow Sz,y\rightarrow x=z$\newline
Clearly, every $x=z$ provable in the system is such that $x$ really equals $z$. Conversley, assume $x=z$. Hence, by definition, $Sx,y\wedge Sz,y$, thus 
meaning the axiom will allow us to conclude that $x=z$.\newline
(d):\newline Take the two axiom schemes:\newline $x<y\rightarrow x\le y$\newline $x=y\rightarrow x\le y$\newline
Any $x\le y$ provable in the system is obviously such that $x$ really $\le y$. Conversley, assume that $x\le y$. Thus, either $x=y$ or $x<y$. Whatever case 
it is, we can use the corresponding axiom to conclude that $x\le y$.\newline
(e):\newline Take the axiom schemes:\newline $x<y\rightarrow x\ne y$\newline $x\ne y\rightarrow y\ne x$
\newline Again, any $x\ne y$ in the systerm obviously is such that $x$ really is not equal to $y$. Conversley, assume that $x\ne y$. Hence, either $x<y$ or 
$y<x$. In the first case, we would immediately conclude $x\ne y$. In the second, we would conclude $y\ne x$ from the first axiom, and the second axiom would 
let us conclude that $x\ne y$.\newline
(f):\newline Take the following $EFS$:\newline $Sx,y\rightarrow x+1=y
$\newline $x+y=z\rightarrow Sy,v\rightarrow Sz,w\rightarrow x+v=w$\newline
Clearly, any $x+y=z$ provable in the system is such that $x+y$ really equals $z$. On the other hand, assume that $x+y=z$. Hence, $x+1=v$ for $Sx,v$ is 
provable. Then, we can use the second axiom $y-1$ times to conclude that $x+y=x+1+y-1=x+y=z$.\newline 
(g):\newline Consider the following $EFS$ over $D$ that has a predicate representing $x\times y=z$:\newline
$x\times 1=x$\newline $x\times y=z\rightarrow Sy,w\rightarrow z+x=v\rightarrow x\times w=v$\newline
Again, every $x\times y=z$ provable in this system clearly is such that $x\times y=z$. Now, assume that $x\times y=z$. Thus, $x\times 1=x$ is provable. Then,
 we apply the second axiom $y-1$ times to recieve $x\times y-1+x=x\times y=z$\newline
(h):\newline We consider the following $EFS$:\newline
$x^1=x$\newline $x^y=z\rightarrow Sy,v\rightarrow x^v=z\times x$\newline	
Every $x^y=z$ provable in the system obviously really means $x^y=z$. On the other hand, assume that $x^y=z$. Thus, $x^1=x$ is provable. Then, we can apply 
the second axiom $y-1$ times to receive $x^y=x\times x^{y-1}=x^y=z$.\newline
\section*{Problem 5}
We consider some dyadically enumerable functional (single-valued) relation R(x,y). We are to show that it is also dyadically solvable. Well, we consider the 
$EFS\ E$ under which $R(x,y)$ is representable by some predicate $P$. Well, we have already shown the relation $x\ne y$ to be dyadically enumerable, and 
therefore, by Problem 1, we consider the $EFS\ E_1$ under which both $R(x,y), x\ne y$ are representable (and $R(x,y)$ is representable by $Px,y$). Well, we 
add the axiom:\newline $Px,y\rightarrow z\ne y\rightarrow Nx,y$\newline Well, because we know that $Px,y$ does not hold for any number other than $y$, 
$Nx,y$ represents $\oldneg R(x,y)$, and hence proves that $R(x,y)$ is dyadically solvable. The proof for trinary relations is similar, replacing $R(x,y)$ by 
$R(x,y,z)$.
\section*{Problem 6}
(a), (c), (f), (g) and (h) are all dyadically solvable by Problem 5 (as they are all functional relations). (b) is dyadically solvable, as $\oldneg (x<y)$ 
can be deterimined by the axioms:\newline $x=y\rightarrow \oldneg (x<y)$\newline $y<x\rightarrow \oldneg (x<y)$, and these can be used as the other relations 
used in these axioms have been shown to be dyadically enumerable. (d) is dyadically solvable, as $\oldneg (x\le y)$ can be determined by:\newline
$y<x\rightarrow \oldneg (x\le y)$\newline (e) is dyadically solvable, as $\oldneg (x\ne y)$ can be given by $x=y$\newline 
\section*{Problem 7}
If $x\in G_n$, it follows that $x$ is either some $g_i\ i<n$ or is a concatenation of strings already in the set; in either case, $1$ must begin $x$. Also, 
$2$ must end $x$ in a similar fashion. Then, at no point can $11Px$, as $x$ is a concatenation of $g_i$'s. Lastly, it is necessary that for $x$ to be in
 $G_n$, $g_{n+1}$ is not a part of $x$ (and this implies $\oldneg (g_mPx)\ m>n$). Hence, we have the following $\Sigma_0$ formula expressing $G_n$:\newline
$x\in G_n$ iff $1Bx\wedge 2Ex\wedge \oldneg (11Px)\wedge \oldneg (g_{n+1}Px)$.
\section*{Problem 8}
Consider the following $\Sigma_0$ formula:\newline
$(\exists v_1\le x)...(\exists v_k\le x)(G_n(v_1)\wedge ...\wedge G_n(v_k)\wedge x=x')$\newline
where $k$ is the number of variables found in the string made up of symbols of $L$ and variables $x_1,...,x_k$ that defines the set $I(X)$, $n$ is from 
the sub-alphabet $k_1,...,k_n$, and $x'$ is the concatenation resulting from, in the formula defining $I(X)$,
 replacing all the occurrences of symbols of $L$ by their respective Gödel numbers, and replacing all occurrences of variables $x_1,...,x_k$ by the 
variables $v_1,...,v_k$ defined in the above formula, respectively.
It is clear that, for any given set $I(X)$, the above formula holds iff $x\in I_0(X)$.   
\section*{Problem 9}
(a):\newline
In general, for a given $L$ with $m-1$ symbols, $x$ is a sequence number of $some$ sequence iff $G_m(x)\wedge g_mBx\wedge g_mEx\wedge \oldneg g_m*g_mPx$
\newline For suppose the formula holds. Then $x$ is a sequence number for $x$, as the $g_m$'s in $x$ are correctly spaced, and for any term other than 
$g_m$ in $x$, they must have a Gödel number less than $g_m$ and are hence $some$ member of $L$. 
 Conversley, assume that $x$ is a sequence number for 
$L$. By defintion, $g_mBx\wedge g_mEx\wedge \oldneg g_m*g_mPx$,
 Also, if $G_m(x)$ did not hold, then there would be some $g_n\ n>m$ in $x$, which is impossible, as there are only $m-1$ terms in $L$. A more specific, 
and completely useless (as it is a restatement of the definition) formula for $x$ being a sequence number of $L$ for some sequence $g_m*X_1*g_m*...*g_m*X_r*
g_m$ is $x=g_m*a_1*g_m*...*g_m*a_r*g_m$.
\newline (b):\newline
$x\in y$ iff $Seq(y)\wedge xPy\wedge \oldneg (g_mPx)$ where $m=n+1$ for $n$ the number of symbols in $L$. $Note:$ Smullyan's answer in the book 
is just another way of saying the above, as it is easy to see that the formula with $G_n(x)$ is equivalent to the one above.\newline
(c):\newline $pr(x,y,z)$ iff $y\in z\wedge (\exists w\le z)(wBz\wedge x\in W\wedge\oldneg (y\in w))$
\section*{Problem 10}
Let us recall that we are considering some $L$ in which the implication sign is part of the ordered alphabet, and therefore has a Gödel number. Thus, 
let $e$ be the Gödel number of the implication sign. Hence, we have the following formula:\newline $Der(x,y,z)$ iff $y=xez\wedge \oldneg (ePx)$ and this 
last condition is necessary as $X$ must be atomic.
\section*{Problem 11}
(a):\newline $Pf(x)$ iff $Seq(x)\wedge (\forall y\le x)[y\in x\supset \newline
((\exists v\le x)(\exists w \le x)(Der(v,w,x)\wedge (pr(w,y,x)\wedge (pr(v,w,x))))\vee A(y))])$ where $A(y)$ is the property of $y$ being the Gödel number 
of an instance of an axiom of $(E)$; and this is $\Sigma_0$ by the Substitution Lemma.\newline
(b):\newline $ypfx$ iff:\newline $Pf(y)\wedge x\in y\wedge (\forall z\le y)(\oldneg pr(z,x,y))$\newline Alternatively (as Smullyan gives in his answer), 
the last condition of the above formula is equivalent to $g_m*x*g_mEy$.\newline
(c):\newline $Pr(x)$ iff $\exists y(ypfx)$\newline
(d):\newline We consider some formally representable set $W_0$ of Gödel numbers of the representable set $W$. So let us consider the predicate $H$ 
that represetns $W_0$, and let its Gödel number be $h$. From this, we have the followng $\Sigma_1$ formula expressing $W_0$:\newline
$\exists y(ypfhx)$\newline
For $n$-ary relations, the above formula can be rewritten as:\newline
$\exists y(ypfha_1c...ca_n)$ where $c$ is the Gödel number of the comma.
\section*{Problem 12}
Suppose $T$ was formally representable. Hence, by Problem 11 (d), it would be expressed by some $\Sigma_1$ formula, thus meaning that it is arithmetic, 
contrary to Tarski's Theorem.
\section*{Problem 13}
Suppose $S$ is a formal arithmetic system that is $\omega$-consistent and is $\Sigma_0$ complete. Well, this means $P_0$ is $\Sigma_1$, and this means 
$P^*$ ($P_0^\#$) is also $\Sigma_1$ (by Problem 12 of Chapter 11), that is, of the form $\exists zR(x,z)$ for some $\Sigma_0$ relation 
$R(x,z)$. Well, because the system is $\Sigma_0$ complete, there is some formula $F(x,z)$ that is provable if $R$ holds, and is 
refutable otherwise. Hence, if $x\in P^*$, $F(x,z)$ holds for at least one $z$, and if $x\notin P^*$, $F(x,z)$ is refutable for every $z$. Thus, $F(x,z)$ 
enumerates $P^*$ (and this result can be generalized for any set). 
Then, if $S$ is a formal system, $P$ is formally representable. Well, consider the $(E)$ in which $P$ is represented by some predicate 
$P$. Then, add the axiom:\newline $Pnx\rightarrow Rx$ for $n$ the Gödel number of the negation, and for some new predicate $R$. Hence $R$ represents $R$, 
and therefore $R$ is also formally representable. Hence both $P^*$ and $R^*$ are enumerable. Then the rest of the problem follows from identical reasoning 
as to that in Chapter 10. 
\end{document}
